% Original analytic teaching diagram, 2026-11-12.
% Omega = [0,1) with length. Each of the 8 cells [m/8,(m+1)/8) shows the binary digits D1, D2, D3 of its points.
% Bottom row: the event {D1 = 1, D3 = 0} = [1/2,5/8) U [3/4,7/8), length 1/4 = (1/2)(1/2).
\definecolor{VB}{HTML}{7188A3}\definecolor{VBF}{HTML}{E5EAF0}
\definecolor{VG}{HTML}{829C8B}\definecolor{VGF}{HTML}{E8EEE9}
\definecolor{VY}{HTML}{B4A16F}\definecolor{VYF}{HTML}{F1ECDD}
\definecolor{VR}{HTML}{A7787D}\definecolor{VRF}{HTML}{F0E5E6}
\definecolor{VP}{HTML}{9588A5}\definecolor{VPF}{HTML}{ECE8F0}
\definecolor{VI}{HTML}{3B4252}\definecolor{VM}{HTML}{6B7280}\definecolor{VL}{HTML}{4C4D4F}
\begin{tikzpicture}[x=1mm,y=1mm,>=Stealth,font=\figurefont\fontsize{8}{10}\selectfont,
lab/.style={text=VM,font=\figurefont\fontsize{7.5}{9}\selectfont},
num/.style={text=VI,font=\figurefont\fontsize{7}{8}\selectfont}]
\def\L{80}
\foreach \k/\y in {1/27,2/20,3/13} {
  \foreach \m in {0,...,7} {
    \pgfmathtruncatemacro{\d}{mod(floor(\m/2^(3-\k)),2)}
    \ifnum\d=1 \fill[VBF] ({\L*\m/8},\y) rectangle ({\L*(\m+1)/8},{\y+5}); \fi
    \draw[VL,line width=.45pt] ({\L*\m/8},\y) rectangle ({\L*(\m+1)/8},{\y+5});
    \node[num] at ({\L*(\m+.5)/8},{\y+2.5}) {\d};
  }
  \node[anchor=east,text=VI] at (-2,{\y+2.5}) {$D_{\k}$};
}
% event row
\draw[VL,line width=.45pt] (0,4) rectangle (\L,9);
\fill[VBF] ({\L*4/8},4) rectangle ({\L*5/8},9); \fill[VBF] ({\L*6/8},4) rectangle ({\L*7/8},9);
\draw[VB,line width=.9pt] ({\L*4/8},4) rectangle ({\L*5/8},9); \draw[VB,line width=.9pt] ({\L*6/8},4) rectangle ({\L*7/8},9);
\node[anchor=east,text=VI] at (-2,6.5) {$\{D_1=1,D_3=0\}$};
\draw[VL,line width=.6pt] (0,0) -- (\L,0);
\foreach \m in {0,...,8} \draw[VL,line width=.6pt] ({\L*\m/8},0) -- ({\L*\m/8},1);
\foreach \x/\t in {0/0,0.5/{1/2},1/1} \node[below,text=VI] at ({\L*\x},-.3) {$\t$};
\node[anchor=west,text=VI] at ({\L+2},0) {$\omega$};
\end{tikzpicture}
